\documentclass[12pt,a4paper]{article}

% --- Настройка языка и шрифтов для LuaLaTeX ---
\usepackage{fontspec}
\defaultfontfeatures{Ligatures=TeX,Scale=MatchLowercase}

\setmainfont{Alice-Regular}[
    Path=font/,
    Extension=.ttf,
    BoldFont = Alice-Regular,
    ItalicFont = Alice-Regular,
    BoldItalicFont = Alice-Regular,
    BoldFeatures = {FakeBold=1.6},
    ItalicFeatures = {FakeSlant=0.2},
    BoldItalicFeatures = {FakeBold=1.6,FakeSlant=0.2}
]

\newfontfamily\cyrillicfont{Alice-Regular}[
    Path=font/,
    Extension=.ttf,
    BoldFont = Alice-Regular,
    ItalicFont = Alice-Regular,
    BoldItalicFont = Alice-Regular,
    BoldFeatures = {FakeBold=1.6},
    ItalicFeatures = {FakeSlant=0.2},
    BoldItalicFeatures = {FakeBold=1.6,FakeSlant=0.2}
]

\usepackage{polyglossia}
\setmainlanguage{russian}
\setsansfont{TeX Gyre Heros}
\setmonofont{TeX Gyre Cursor}

\usepackage{graphicx}        					          % Графика
\usepackage[protrusion=false]{microtype}       	% Микротипографика
\usepackage{amsmath, amsfonts, amssymb, amsthm}	% Математика
\usepackage{braket} 							              % Braket нотация для квантовой механики
\usepackage{unicode-math} 						          % Современный пакет для математики в LuaLaTeX
\usepackage{longtable}       					          % Таблицы с разрывом страниц
\usepackage{tikz}           					          % Рисование
\usepackage[europeanresistors]{circuitikz}		  % Рисование электрических схем в европейском стиле
\usepackage{listings}							              % Для отображения кода
\usepackage{tabularx}       					          % Для таблиц во всю ширину
\usepackage{array}		  						            % Для расширенных возможностей таблиц
\usepackage{float}          					          % Для точного позиционирования фигур
\usepackage{pgfplots}                           % Для построения графиков
\usepackage{caption}                            % Для настройки подписей к рисункам и таблицам
\usepackage{xfp}                                % Для вычислений в LaTeX
\pgfplotsset{compat=1.18}
\usepackage[a4paper,left=20mm,right=20mm,top=20mm,bottom=20mm]{geometry}

\usetikzlibrary{arrows.meta}

% --- Колонтитулы ---
\usepackage{fancyhdr}
\pagestyle{fancy}
\fancyhf{}

\renewcommand{\headrulewidth}{0.4pt}
\renewcommand{\footrulewidth}{0.4pt}
\renewcommand\lstlistingname{Исходный код}

\lstdefinestyle{main}{
	backgroundcolor=\color{gray!8},   
	commentstyle=\color{green!50!black},
	keywordstyle=\color{blue!80!black},
	numberstyle=\small\color{darkgray},
	stringstyle=\rmfamily\color{orange!90!black},
	basicstyle=\ttfamily\small,
	breakatwhitespace=false,         
	breaklines=true,                   
	keepspaces=true,                 
	numbers=left,                    
	numbersep=8pt,                  
	showspaces=false,                
	showstringspaces=false,
	showtabs=false,                  
	tabsize=4
}	

\lstset{style=main}

\newcolumntype{C}{>{\centering\arraybackslash}X}
\renewcommand{\arraystretch}{1.25}

% =========================================================
%         СЛУЖЕБНЫЕ КОМАНДЫ ДЛЯ ВЫЧИСЛЕНИЙ
% =========================================================
\newcommand{\CalcZero} [1]{\fpeval{round((#1),0)}}
\newcommand{\CalcOne}  [1]{\fpeval{round((#1),1)}}
\newcommand{\CalcTwo}  [1]{\fpeval{round((#1),2)}}
\newcommand{\CalcThree}[1]{\fpeval{round((#1),3)}}
\newcommand{\CalcFour} [1]{\fpeval{round((#1),4)}}
\newcommand{\CalcFive} [1]{\fpeval{round((#1),5)}}

% ---------------------------------------------------------
% ОБЩИЕ ДАННЫЕ
% ---------------------------------------------------------
\def\VariantN{1}          					         % Номер варианта
\def\StudentName{Андреев Иван Алексеевич}  	 % ФИО студента
\def\LabNumber{15}							               % Номер лабораторной работы
\def\Discipline{Атомная физика}				           % Название дисциплины
\def\LabTitle{Волны де Бройля и дифракция электронов}	 % Название лабораторной работы

\def\dH{\fpeval{68 * 10^(-3)}} % Межслойное расстояние в кристалле, м
\def\D {\fpeval{127 * 10^(-3)}} % Расстояние от кристалла до экрана, м
\def\h {\fpeval{6.62607015 * 10^(-34)}} % Постоянная Планка, Дж*с
\def\me{\fpeval{9.10938356 * 10^(-31)}} % Масса электрона, кг
\def\ee{\fpeval{1.602176634 * 10^(-19)}} % Заряд электрона, Кл 
\def\c {\fpeval{299792458}} % Скорость света, м/с

\def\UzeroOne            {\fpeval{3.5 * 10^3}} % Напряжение, В
\def\UzeroTwo            {\fpeval{4.0 * 10^3}} % Напряжение, В
\def\UzeroThree          {\fpeval{4.5 * 10^3}} % Напряжение, В
\def\UzeroFour           {\fpeval{5.0 * 10^3}} % Напряжение, В
\def\AdditiveConstant    {\fpeval{1.7 * 10^(-2)}} % Константа измерений, м
\def\rOneThreePointFivekV{\fpeval{((0.5 + 1.3) / 2 * 10^(-2) + \AdditiveConstant) / 2}} % Радиус первой окружности, м для 3.5 кВ
\def\rTwoThreePointFivekV{\fpeval{((2.6 + 3.5) / 2 * 10^(-2) + \AdditiveConstant) / 2}} % Радиус второй окружности, м для 3.5 кВ
\def\rOneFourPointZerokV {\fpeval{((0.6 + 1.2) / 2 * 10^(-2) + \AdditiveConstant) / 2}} % Радиус первой окружности, м для 4.0 кВ
\def\rTwoFourPointZerokV {\fpeval{((2.3 + 3.0) / 2 * 10^(-2) + \AdditiveConstant) / 2}} % Радиус второй окружности, м для 4.0 кВ
\def\rOneFourPointFivekV {\fpeval{((0.5 + 1.0) / 2 * 10^(-2) + \AdditiveConstant) / 2}} % Радиус первой окружности, м для 4.5 кВ
\def\rTwoFourPointFivekV {\fpeval{((2.0 + 2.8) / 2 * 10^(-2) + \AdditiveConstant) / 2}} % Радиус второй окружности, м для 4.5 кВ
\def\rOneFivePointZerokV {\fpeval{((0.4 + 0.9) / 2 * 10^(-2) + \AdditiveConstant) / 2}} % Радиус первой окружности, м для 5.0 кВ
\def\rTwoFivePointZerokV {\fpeval{((2.0 + 2.5) / 2 * 10^(-2) + \AdditiveConstant) / 2}} % Радиус второй окружности, м для 5.0 кВ



\title{
	МИНОБРНАУКИ РОССИИ

	федеральное государственное автономное образовательное учреждение

	высшего образования

	«Национальный исследовательский университет
	«Московский институт электронной техники»
	\vspace{2em}
}
\author{
}
\date{}

\begin{document}

\maketitle
\thispagestyle{fancy}

\begin{center}
\Large Отчёт по Лабораторной работе $\LabNumber$\\
по дисциплине «\Discipline»\\
«\LabTitle»
\end{center}

\begin{center}
\Large Вариант \VariantN
\end{center}

\vspace{12em}

\cfoot{Москва 2026}

\pagebreak

\fancyhead[L]{\LabTitle} 
\fancyhead[R]{\Discipline}
\fancyfoot[C]{\thepage}

\begin{tabularx}{\textwidth}{|C|C|C|C|C|C|}
  \hline
  $N$ п/п & $U_0$, В & $2r_1$, мм & $2r_2$, мм & $1/r_1^2$, м$^{-2}$ & $1/r_2^2$, м$^{-2}$
  \tabularnewline
  \hline
  1 & $\UzeroOne$ & $\CalcFive{\rOneThreePointFivekV * 10^3 * 2}$ & 
  $\CalcFive{\rTwoThreePointFivekV * 10^3 * 2}$ & 
  $\CalcFive{1/(\rOneThreePointFivekV^2)}$ & 
  $\CalcFive{1/(\rTwoThreePointFivekV^2)}$ \tabularnewline
  2 & $\UzeroTwo$ & $\CalcFive{\rOneFourPointZerokV * 10^3 * 2}$ & 
  $\CalcFive{\rTwoFourPointZerokV * 10^3 * 2}$ & 
  $\CalcFive{1/(\rOneFourPointZerokV^2)}$ & 
  $\CalcFive{1/(\rTwoFourPointZerokV^2)}$ \tabularnewline
  3 & $\UzeroThree$ & $\CalcFive{\rOneFourPointFivekV * 10^3 * 2}$ & 
  $\CalcFive{\rTwoFourPointFivekV * 10^3 * 2}$ & 
  $\CalcFive{1/(\rOneFourPointFivekV^2)}$ & 
  $\CalcFive{1/(\rTwoFourPointFivekV^2)}$ \tabularnewline
  4 & $\UzeroFour$ & $\CalcFive{\rOneFivePointZerokV * 10^3 * 2}$ & 
  $\CalcFive{\rTwoFivePointZerokV * 10^3 * 2}$ & 
  $\CalcFive{1/(\rOneFivePointZerokV^2)}$ & 
  $\CalcFive{1/(\rTwoFivePointZerokV^2)}$ \tabularnewline
  \hline
\end{tabularx}

% Абсолютная погрешность измерения радиуса
\def\DeltaR{\fpeval{0.05 * 10^(-2)}} % м
\def\DeltaU{\fpeval{0.1 * 10^3}} % В

% Значения 1/r^2
\def\YOneThreePointFivekV{\fpeval{1 / \rOneThreePointFivekV^2}}
\def\YOneFourPointZerokV {\fpeval{1 / \rOneFourPointZerokV^2}}
\def\YOneFourPointFivekV {\fpeval{1 / \rOneFourPointFivekV^2}}
\def\YOneFivePointZerokV {\fpeval{1 / \rOneFivePointZerokV^2}}

\def\YTwoThreePointFivekV{\fpeval{1 / \rTwoThreePointFivekV^2}}
\def\YTwoFourPointZerokV {\fpeval{1 / \rTwoFourPointZerokV^2}}
\def\YTwoFourPointFivekV {\fpeval{1 / \rTwoFourPointFivekV^2}}
\def\YTwoFivePointZerokV {\fpeval{1 / \rTwoFivePointZerokV^2}}

% Погрешности для 1/r^2:
% Δ(1/r^2) = 2 Δr / r^3
\def\DeltaYOneThreePointFivekV{\fpeval{2 * \DeltaR / \rOneThreePointFivekV^3}}
\def\DeltaYOneFourPointZerokV {\fpeval{2 * \DeltaR / \rOneFourPointZerokV^3}}
\def\DeltaYOneFourPointFivekV {\fpeval{2 * \DeltaR / \rOneFourPointFivekV^3}}
\def\DeltaYOneFivePointZerokV {\fpeval{2 * \DeltaR / \rOneFivePointZerokV^3}}

\def\DeltaYTwoThreePointFivekV{\fpeval{2 * \DeltaR / \rTwoThreePointFivekV^3}}
\def\DeltaYTwoFourPointZerokV {\fpeval{2 * \DeltaR / \rTwoFourPointZerokV^3}}
\def\DeltaYTwoFourPointFivekV {\fpeval{2 * \DeltaR / \rTwoFourPointFivekV^3}}
\def\DeltaYTwoFivePointZerokV {\fpeval{2 * \DeltaR / \rTwoFivePointZerokV^3}}

% Коэффициенты прямых, полученные по линейной аппроксимации:
% 1/r^2 = gamma * (U0 - U_R)
\def\GammaApproxOne{\fpeval{0.9449}}
\def\URapproxOne    {\fpeval{2.4195 * 10^3}}

\def\GammaMaxOne{((\YOneThreePointFivekV - \DeltaYOneThreePointFivekV) - (\YOneFivePointZerokV + \DeltaYOneFivePointZerokV)) / ((\UzeroOne + \DeltaU) - (\UzeroFour - \DeltaU))}
\def\URMaxOne   {(\YOneFivePointZerokV + \DeltaYOneFivePointZerokV) + \GammaMaxOne * -(\UzeroFour - \DeltaU)}

\def\GammaMinOne{((\YOneFourPointZerokV + \DeltaYOneFourPointZerokV) - (\YOneFivePointZerokV - \DeltaYOneFivePointZerokV)) / ((\UzeroTwo -  \DeltaU) - (\UzeroFour + \DeltaU))}
\def\URMinOne   {(\YOneFivePointZerokV - \DeltaYOneFivePointZerokV) + \GammaMinOne * -(\UzeroFour + \DeltaU)}

\def\GammaApproxTwo{\fpeval{0.5276}}
\def\URapproxTwo    {\fpeval{-34.9495}}

\def\GammaMaxTwo{((\YTwoThreePointFivekV - \DeltaYTwoThreePointFivekV) - (\YTwoFivePointZerokV + \DeltaYTwoFivePointZerokV)) / ((\UzeroOne + \DeltaU) - (\UzeroFour - \DeltaU))}
\def\URMaxTwo   {(\YTwoFivePointZerokV + \DeltaYTwoFivePointZerokV) + \GammaMaxTwo * -(\UzeroFour - \DeltaU)}

\def\GammaMinTwo{((\YTwoThreePointFivekV + \DeltaYTwoThreePointFivekV) - (\YTwoFivePointZerokV - \DeltaYTwoFivePointZerokV)) / ((\UzeroOne -  \DeltaU) - (\UzeroFour + \DeltaU))}
\def\URMinTwo   {(\YTwoFivePointZerokV - \DeltaYTwoFivePointZerokV) + \GammaMinTwo * -(\UzeroFour + \DeltaU)}

\begin{equation*}
    \begin{aligned}
        \frac{1}{r^2} \quad & \left(\frac{1}{r_1^2}\right)_1 = \YOneThreePointFivekV \quad & \left(\frac{1}{r_2^2}\right)_1 = \YTwoThreePointFivekV \\
        & \left(\frac{1}{r_1^2}\right)_2 = \YOneFourPointZerokV \quad & \left(\frac{1}{r_2^2}\right)_2 = \YTwoFourPointZerokV \\
        & \left(\frac{1}{r_1^2}\right)_3 = \YOneFourPointFivekV \quad & \left(\frac{1}{r_2^2}\right)_3 = \YTwoFourPointFivekV \\
        & \left(\frac{1}{r_1^2}\right)_4 = \YOneFivePointZerokV \quad & \left(\frac{1}{r_2^2}\right)_4 = \YTwoFivePointZerokV
    \end{aligned}
\end{equation*}

\begin{equation*}
    \begin{aligned}
        \Delta\left(\frac{1}{r^2}\right) = \frac{2 \Delta r}{r^3} 
        \quad & \Delta\left(\frac{1}{r_1^2}\right)_1 = \frac{2 \cdot \DeltaR}{\rOneThreePointFivekV^3} = \DeltaYOneThreePointFivekV \\
        & \Delta\left(\frac{1}{r_1^2}\right)_2 = \frac{2 \cdot \DeltaR}{\rOneFourPointZerokV^3} = \DeltaYOneFourPointZerokV \\
        & \Delta\left(\frac{1}{r_1^2}\right)_3 = \frac{2 \cdot \DeltaR}{\rOneFourPointFivekV^3} = \DeltaYOneFourPointFivekV \\
        & \Delta\left(\frac{1}{r_1^2}\right)_4 = \frac{2 \cdot \DeltaR}{\rOneFivePointZerokV^3} = \DeltaYOneFivePointZerokV \\
        & \Delta\left(\frac{1}{r_2^2}\right)_1 = \frac{2 \cdot \DeltaR}{\rTwoThreePointFivekV^3} = \DeltaYTwoThreePointFivekV \\
        & \Delta\left(\frac{1}{r_2^2}\right)_2 = \frac{2 \cdot \DeltaR}{\rTwoFourPointZerokV^3} = \DeltaYTwoFourPointZerokV \\
        & \Delta\left(\frac{1}{r_2^2}\right)_3 = \frac{2 \cdot \DeltaR}{\rTwoFourPointFivekV^3} = \DeltaYTwoFourPointFivekV \\
        & \Delta\left(\frac{1}{r_2^2}\right)_4 = \frac{2 \cdot \DeltaR}{\rTwoFivePointZerokV^3} = \DeltaYTwoFivePointZerokV
    \end{aligned}
\end{equation*}

\begin{equation*}
    \begin{matrix}
        & ^1\gamma_\text{approx} = \GammaApproxOne 
        & ^1\gamma_\text{max} = \CalcFive{\GammaMaxOne}
        & ^1\gamma_\text{min} = \CalcFive{\GammaMinOne} \\
        & ^1U_R^\text{approx} = \URapproxOne 
        & ^1U_R^\text{max} = \CalcFive{\URMaxOne} 
        & ^1U_R^\text{min} = \CalcFive{\URMinOne} \\ \\
        & ^2\gamma_\text{approx} = \GammaApproxTwo 
        & ^2\gamma_\text{max} = \CalcFive{\GammaMaxTwo}
        & ^2\gamma_\text{min} = \CalcFive{\GammaMinTwo} \\
        & ^2U_R^\text{approx} = \URapproxTwo 
        & ^2U_R^\text{max} = \CalcFive{\URMaxTwo} 
        & ^2U_R^\text{min} = \CalcFive{\URMinTwo}
    \end{matrix}
\end{equation*}

\begin{figure}[H]
\centering
\begin{tikzpicture}
\begin{axis}[
    width=\textwidth,
    height=\textwidth,
    title={Зависимость $\frac{1}{r^2}$ от ускоряющего напряжения $U_0$},
    xlabel={$U_0$, В},
    ylabel={$\frac{1}{r^2}$, м$^{-2}$},
    xmin=\fpeval{min(\UzeroOne, \UzeroTwo, \UzeroThree, \UzeroFour) * 0.9},
    xmax=\fpeval{max(\UzeroOne, \UzeroTwo, \UzeroThree, \UzeroFour) * 1.1},
    ymin=\fpeval{min(\YOneThreePointFivekV, \YOneFourPointZerokV, \YOneFourPointFivekV, \YOneFivePointZerokV, 
                     \YTwoThreePointFivekV, \YTwoFourPointZerokV, \YTwoFourPointFivekV, \YTwoFivePointZerokV) * 0.9},
    ymax=\fpeval{max(\YOneThreePointFivekV, \YOneFourPointZerokV, \YOneFourPointFivekV, \YOneFivePointZerokV, 
                     \YTwoThreePointFivekV, \YTwoFourPointZerokV, \YTwoFourPointFivekV, \YTwoFivePointZerokV) * 1.1},
    xtick={},
    ytick={},
    minor x tick num=9,
    minor y tick num=9,
    grid=both,
    major grid style={
        line width=0.45pt,
        draw=black!55
    },
    minor grid style={
        line width=0.15pt,
        draw=black!20
    },
    tick align=outside,
    axis line style={black},
    tick style={black},
    legend style={
        at={(0.03,0.97)},
        anchor=north west,
        draw=black,
        fill=white,
        font=\small
    },
    every axis plot/.append style={line width=1pt},
]

% ---------------- Первое кольцо ----------------

\addplot+[
    only marks,
    mark=*,
    mark size=2pt,
    error bars/.cd,
        y dir=both,
        y explicit,
        x dir=both,
        x explicit
]
coordinates {
    (\UzeroOne,   \YOneThreePointFivekV) +- (\DeltaU,\DeltaYOneThreePointFivekV)
    (\UzeroTwo,   \YOneFourPointZerokV)  +- (\DeltaU,\DeltaYOneFourPointZerokV)
    (\UzeroThree, \YOneFourPointFivekV)  +- (\DeltaU,\DeltaYOneFourPointFivekV)
    (\UzeroFour,  \YOneFivePointZerokV)  +- (\DeltaU,\DeltaYOneFivePointZerokV)
};

\addlegendentry{Первое кольцо}

\addplot[
    domain=3400:5100,
    samples=2,
    dashed,
    color=green
]
{\GammaApproxOne * x + \URapproxOne};

\addlegendentry{Аппроксимация $r_1$}

% ---------------- Второе кольцо ----------------

\addplot+[
    only marks,
    mark=*,
    mark size=2pt,
    error bars/.cd,
        y dir=both,
        y explicit,
        x dir=both,
        x explicit
]
coordinates {
    (\UzeroOne,   \YTwoThreePointFivekV) +- (\DeltaU,\DeltaYTwoThreePointFivekV)
    (\UzeroTwo,   \YTwoFourPointZerokV)  +- (\DeltaU,\DeltaYTwoFourPointZerokV)
    (\UzeroThree, \YTwoFourPointFivekV)  +- (\DeltaU,\DeltaYTwoFourPointFivekV)
    (\UzeroFour,  \YTwoFivePointZerokV)  +- (\DeltaU,\DeltaYTwoFivePointZerokV)
};

\addlegendentry{Второе кольцо}

\addplot[
    domain=3400:5100,
    samples=2,
    dashed,
    color=green
]
{\GammaApproxTwo * x + \URapproxTwo};

\addlegendentry{Аппроксимация $r_2$}

\addplot[
    domain=3400:5100,
    samples=2,
    dashed,
    color=red
]
{(\GammaMaxOne * x + \URMaxOne)};

\addplot[
    domain=3400:5100,
    samples=2,
    dashed,
    color=red
]
{(\GammaMinOne * x + \URMinOne)};

\addplot[
    domain=3400:5100,
    samples=2,
    dashed,
    color=red
]
{(\GammaMaxTwo * x + \URMaxTwo)};

\addplot[
    domain=3400:5100,
    samples=2,
    dashed,
    color=red
]
{(\GammaMinTwo * x + \URMinTwo)};

\end{axis}
\end{tikzpicture}

\caption{Графики зависимости $\frac{1}{r^2}(U_0)$ для двух первых дифракционных колец}
\label{fig:de-broglie-r2-u0}
\end{figure}

\begin{equation*}
    \frac{1}{r^2} = \frac{2m_eed^2}{D^2h^2}(U_0 - U_R) \quad d_\text{н} = 68 \text{ мм} \quad D = 127 \text{ мм}
\end{equation*}

\def\GammaMeanOne{\fpeval{(\GammaMaxOne + \GammaMinOne) / 2}}
\def\GammaMeanTwo{\fpeval{(\GammaMaxTwo + \GammaMinTwo) / 2}}
\def\GammaDeltaOne{\fpeval{(\GammaMaxOne - \GammaMinOne) / 2}}
\def\GammaDeltaTwo{\fpeval{(\GammaMaxTwo - \GammaMinTwo) / 2}}

\begin{equation*}
    \begin{aligned}
        \braket{\gamma} = \frac{\gamma_\text{max} + \gamma_\text{min}}{2} \quad & \braket{\gamma}_1 = \frac{\CalcFive{\GammaMaxOne} + \CalcFive{\GammaMinOne}}{2} = \CalcFive{\GammaMeanOne} \\
        & \braket{\gamma}_2 = \frac{\CalcFive{\GammaMaxTwo} + \CalcFive{\GammaMinTwo}}{2} = \CalcFive{\GammaMeanTwo}
    \end{aligned}
\end{equation*}

\begin{equation*}
    \begin{aligned}
        \braket{\Delta\gamma} = \frac{\gamma_\text{max} - \gamma_\text{min}}{2} \quad & \braket{\Delta\gamma}_1 = \frac{\CalcFive{\GammaMaxOne} - \CalcFive{\GammaMinOne}}{2} = \CalcFive{\GammaDeltaOne} \\
        & \braket{\Delta\gamma}_2 = \frac{\CalcFive{\GammaMaxTwo} - \CalcFive{\GammaMinTwo}}{2} = \CalcFive{\GammaDeltaTwo}
    \end{aligned}
\end{equation*}

\def\dOne{\fpeval{\D *  * sqrt(\GammaMeanOne / (2 * \me * \ee))}}
\def\dTwo{\fpeval{\D * \h * sqrt(\GammaMeanTwo / (2 * \me * \ee))}}

\begin{equation*}
    \begin{aligned}
        \gamma = \frac{2m_eed^2}{D^2h^2} \quad & d_1 = \CalcFive{\D} \cdot h \cdot \sqrt{\frac{\CalcFive{\GammaMeanOne}}{2 \cdot m_e \cdot e}} = \CalcFive{\dOne * 10^9} \cdot 10^{-9} \\
        \quad d = Dh \sqrt{\frac{\gamma}{2m_ee}} \quad & d_2 = \CalcFive{\D} \cdot h \cdot \sqrt{\frac{\CalcFive{\GammaMeanTwo}}{2 \cdot m_e \cdot e}} = \CalcFive{\dTwo * 10^9} \cdot 10^{-9}
    \end{aligned}
\end{equation*}

\def\DeltaDOne{\fpeval{\D * \h * \GammaDeltaOne / (2 * sqrt(2 * \me * \ee * \GammaMeanOne))}}
\def\DeltaDTwo{\fpeval{\D * \h * \GammaDeltaTwo / (2 * sqrt(2 * \me * \ee * \GammaMeanTwo))}}

\begin{equation*}
    \begin{aligned}
        \Delta d = Dh \frac{\Delta\gamma}{2\sqrt{2m_ee\gamma}} \quad & \Delta d_1 = \CalcFive{\D} \cdot h \cdot \frac{\CalcFive{\GammaDeltaOne}}{2 \cdot \sqrt{2 \cdot m_e \cdot e \cdot \CalcFive{\GammaMeanOne}}} = \CalcFive{\DeltaDOne * 10^9} \cdot 10^{-9} \\
        & \Delta d_2 = \CalcFive{\D} \cdot h \cdot \frac{\CalcFive{\GammaDeltaTwo}}{2 \cdot \sqrt{2 \cdot m_e \cdot e \cdot \CalcFive{\GammaMeanTwo}}} = \CalcFive{\DeltaDTwo * 10^9} \cdot 10^{-9}
    \end{aligned}
\end{equation*}

\begin{figure}[H]
\centering
\begin{tikzpicture}
\begin{axis}[
    width=\textwidth,
    height=\textwidth,
    title={Зависимость $\frac{1}{r^2}$ от ускоряющего напряжения $U_0$},
    xlabel={$U_0$, В},
    ylabel={$\frac{1}{r^2}$, м$^{-2}$},
    xmin=\fpeval{min(\UzeroOne, \UzeroTwo, \UzeroThree, \UzeroFour) * -0.9},
    xmax=\fpeval{max(\UzeroOne, \UzeroTwo, \UzeroThree, \UzeroFour) * 1.1},
    ymin=\fpeval{min(\YOneThreePointFivekV, \YOneFourPointZerokV, \YOneFourPointFivekV, \YOneFivePointZerokV, 
                     \YTwoThreePointFivekV, \YTwoFourPointZerokV, \YTwoFourPointFivekV, \YTwoFivePointZerokV) * -0.1},
    ymax=\fpeval{max(\YOneThreePointFivekV, \YOneFourPointZerokV, \YOneFourPointFivekV, \YOneFivePointZerokV, 
                     \YTwoThreePointFivekV, \YTwoFourPointZerokV, \YTwoFourPointFivekV, \YTwoFivePointZerokV) * 1.1},
    xtick={},
    ytick={},
    minor x tick num=9,
    minor y tick num=9,
    grid=both,
    major grid style={
        line width=0.45pt,
        draw=black!55
    },
    minor grid style={
        line width=0.15pt,
        draw=black!20
    },
    tick align=outside,
    axis line style={black},
    tick style={black},
    legend style={
        at={(0.03,0.97)},
        anchor=north west,
        draw=black,
        fill=white,
        font=\small
    },
    every axis plot/.append style={line width=1pt},
]

% ---------------- Первое кольцо ----------------

\addplot+[
    only marks,
    mark=*,
    mark size=2pt,
    error bars/.cd,
        y dir=both,
        y explicit,
        x dir=both,
        x explicit
]
coordinates {
    (\UzeroOne,   \YOneThreePointFivekV) +- (\DeltaU,\DeltaYOneThreePointFivekV)
    (\UzeroTwo,   \YOneFourPointZerokV)  +- (\DeltaU,\DeltaYOneFourPointZerokV)
    (\UzeroThree, \YOneFourPointFivekV)  +- (\DeltaU,\DeltaYOneFourPointFivekV)
    (\UzeroFour,  \YOneFivePointZerokV)  +- (\DeltaU,\DeltaYOneFivePointZerokV)
};

\addlegendentry{Первое кольцо}

\addplot[
    domain=-3000:5100,
    samples=2,
    dashed,
    color=green
]
{\GammaApproxOne * x + \URapproxOne};

\addlegendentry{Первое значение $b_1 = \URapproxOne$}

% ---------------- Второе кольцо ----------------

\addplot+[
    only marks,
    mark=*,
    mark size=2pt,
    error bars/.cd,
        y dir=both,
        y explicit,
        x dir=both,
        x explicit
]
coordinates {
    (\UzeroOne,   \YTwoThreePointFivekV) +- (\DeltaU,\DeltaYTwoThreePointFivekV)
    (\UzeroTwo,   \YTwoFourPointZerokV)  +- (\DeltaU,\DeltaYTwoFourPointZerokV)
    (\UzeroThree, \YTwoFourPointFivekV)  +- (\DeltaU,\DeltaYTwoFourPointFivekV)
    (\UzeroFour,  \YTwoFivePointZerokV)  +- (\DeltaU,\DeltaYTwoFivePointZerokV)
};

\addlegendentry{Второе кольцо}

\addplot[
    domain=0:5100,
    samples=2,
    dashed,
    color=green
]
{\GammaApproxTwo * x + \URapproxTwo};

\addlegendentry{Второе значение $b_2 = \URapproxTwo$}

\end{axis}
\end{tikzpicture}
\end{figure}

\begin{equation*}
    \begin{aligned}
        U_R = -\frac{b}{\gamma} \quad & U_{R1} = \CalcFive{-\URapproxOne / \GammaApproxOne} \\
        & U_{R2} = \CalcFive{-\URapproxTwo / \GammaApproxTwo}
    \end{aligned}
\end{equation*}

\end{document}